% Solutions to the Chapter 10 exercises that are not code, from lectures/L10/appendix/c_exercises.md.
% Exercises 10.6 and 10.7 are Code and are not answered; see chapter.tex for why.

\section{Chapter 10: The Device Model and the Device Tree}
\label{sol:sec:c10}
Answers to the exercises in \secref{c10:sec:exercises}. Exercises~\ref{c10:ex:6} and~\ref{c10:ex:7}
are Code, and are checked on the target: the first as its Check yourself line says, the second by
the unbind and rebind that \file{run.sh} performs.

% ------------------------------------------------------------------------------------------
\solution[fillaccent2]{c10:ex:1}{Bus, device, driver}{Recall}

\xpart{a} \textbf{A bus} is a matching policy: the rule by which the kernel decides which of the
drivers registered with it can handle which of the devices registered with it, and the two lists it
applies the rule to. \textbf{A device} is the kernel's record that something exists, with whatever
describes it: its identity, its registers, its interrupts. \textbf{A driver} is code that says which
devices it can handle and supplies the \code{probe} and \code{remove} that take one on and give it
up.

\xpart{b} It is for devices that cannot be discovered: the blocks inside an SoC, which sit at fixed
addresses on the processor's own interconnect and announce nothing. That is also what its devices
share: something other than the hardware has to declare that each one exists, and on this target
that something is the device tree. The bus then matches them against drivers by
\code{compatible}, by ACPI identity or by name (\code{platform_match},
\file{drivers/base/platform.c}).

\xpart{c} The kernel created it, from the \code{qa-dev@0} node in the device tree QEMU generated,
during boot and before any module could be loaded. \code{of_platform_default_populate_init}, an
\code{arch_initcall_sync} (\file{drivers/of/platform.c}:600), walks the tree from the root. It
descends into \code{platform-bus@c000000} because that node's \code{compatible} includes
\code{simple-bus}, and creates a platform device for each child. \code{of_device_alloc} translates
the node's \code{reg} into a resource there and then, and the device is named \code{c000000.qa-dev}
from that translated address and the node's name (\code{of_device_make_bus_id},
\file{drivers/of/device.c}:302--309).

\xpart{d} No; either order works, and that is the point of the model. \textbf{Device first}, as
\code{qa-dev} is: registering the driver makes the bus try it against every device not yet bound,
and \code{probe} runs inside \code{insmod}. \textbf{Driver first}: the driver waits, and when the
device is added later, because it was plugged in, or an overlay added its node, or the driver of
its parent probed, the bus tries the new device against every registered driver and \code{probe}
runs then. The author does nothing differently for either: \code{probe} is written the same way.
What the author must not do is assume an order between devices, which is what
\code{-EPROBE_DEFER} is for.

\xpart{e} \file{/sys/bus/<bus>/devices} groups devices by how they are attached, which is the
matching policy that found them; \file{/sys/class/<class>} groups them by what they are to
userspace, whatever they are attached by. A program looking for all the serial ports reads
\file{/sys/class/tty}, where a PL011, a USB serial adapter and a virtual console all appear, although
the first is on the \code{amba} bus, the second on \code{usb}, and the third on no bus at all.

% ------------------------------------------------------------------------------------------
\solution[fillaccent2]{c10:ex:2}{Probe and remove}{Recall}

\xpart{a} Once for each device that matches: two \code{qa-dev} nodes, or two of the same chip on a
board, mean two calls to the same function in the same loaded module. It also runs again for a
device that is unbound and bound again, and after a deferral. That rules out keeping anything that
belongs to a device in a file-scope variable. Everything a device needs hangs off its own
\code{struct device}: allocated with \code{devm_kzalloc} against it, attached with
\code{platform_set_drvdata}, and found again from whatever pointer each callback is handed.

\xpart{b} Any machine with two of the device, as in part e) of Exercise~\ref{c10:ex:9}. The second
\code{probe} overwrites \code{regs}, so from then on everything done for the first device goes to
the second device's registers. The first device's handler reads and acknowledges the second
device's \code{IRQ_STATUS}, never its own, so the first device's level-triggered line is never
cleared: an interrupt storm (\secref{c8:app:a6}), or, if the handler correctly returns
\code{IRQ_NONE} when the status it read shows nothing pending, the kernel disabling the line with
\code{nobody cared} (\secref{c8:app:a:the-return-value}). And unbinding the second device unmaps the
registers the first is still using.

\xpart{c} Because by the time \code{remove} is called, the decision has been made somewhere the
driver has no say: \code{rmmod}, a write to \code{unbind}, or the device itself going away, perhaps
already physically gone. The core has nothing to offer in return for an error. It cannot keep a
driver bound to a device that is being removed, or unload half a module, and a \code{remove} that
stops partway leaves state that nobody can finish tearing down. So a failure from \code{remove}
would mean nothing, and the \code{void} return type (\file{include/linux/platform_device.h}:245)
says so. What the driver must do instead is make every step of \code{remove} succeed, which with
\code{devm_} is little more than stopping the device.

\xpart{d} It says ``something I need is not available yet; call me again later'': a clock, a
regulator, a GPIO or an interrupt controller whose driver has not probed. The core puts the device
on its list of deferred devices without reporting a failure (\code{driver_deferred_probe_add},
\file{drivers/base/dd.c}:132), tries every device on that list again each time another driver binds
successfully, and once more when the initcalls are done; anything still waiting is listed in
\file{/sys/kernel/debug/devices_deferred}. The functions that look up a dependency, such as
\code{devm_clk_get} and \code{platform_get_irq}, return \code{-EPROBE_DEFER} themselves, so a driver
usually only passes it on. The alternative is to return a real error such as \code{-ENODEV}: the
probe fails for good, and the driver works only on a kernel where the dependency happened to probe
first, which is decided by link and initcall order and not by anything the driver controls.

\xpart{e} Whatever makes the device reachable from userspace: registering a character device, a
class device, or the device with a subsystem such as IIO or GPIO. The reason is the one
\secref{c5:app:a3} gave: after \code{cdev_add} the device is live, and a call can arrive in your
code, on another CPU, before \code{probe} has returned. So everything that call touches, the
mapping, the interrupt, the private structure, must be set up before the call that publishes it.
The driver core keeps the same rule for the attributes in \code{.dev_groups}: it adds them only
after \code{probe} has returned successfully (\file{drivers/base/dd.c}:676).

% ------------------------------------------------------------------------------------------
\solution[fillaccent3]{c10:ex:3}{Reading a device tree}{Hand calculation}

\xpart{a} Two: one address cell and one size cell. You know because \code{thing}'s parent,
\code{bus@10000000}, declares \code{\#address-cells = <1>} and \code{\#size-cells = <1>}; the node
itself does not say. So \code{reg = <0x2000 0x40>} is one region, at address \code{0x2000} and
\code{0x40} bytes long.

\xpart{b} Four. The child address takes the bus's own \code{\#address-cells}, 1: \code{0x1000}. The
parent address takes the root's \code{\#address-cells}, 2: \code{0x0 0x10000000}, which is
\code{0x0_10000000}. The size takes the bus's \code{\#size-cells}, 1: \code{0x100000}. So child
addresses \code{0x1000} to \code{0x100FFF} map to parent addresses \code{0x10000000} to
\code{0x100FFFFF}.

\xpart{c} \code{0x2000} lies inside the child range, which starts at \code{0x1000}, so the offset
into the range is \code{0x2000 - 0x1000 = 0x1000}, and the registers start at
\code{0x10000000 + 0x1000 =} \textbf{\code{0x10001000}}. They end at
\code{0x10001000 + 0x40 - 1 = 0x1000103F}. Adding the child address to the parent base without
subtracting the child base, which is the mistake a range that does not start at zero exists to
catch, gives \code{0x10002000}, one page too high.

\xpart{d} \code{acme,thing-v2}. A driver binds if its table contains any of the node's strings, and
where the table contains several, the entry for the earlier string scores higher
(\code{__of_device_is_compatible}, \file{drivers/of/base.c}:334), so a driver that knows both uses
its \code{thing-v2} entry and whatever match data goes with it. The second string is the fallback:
it says the chip is compatible with the original \code{thing}, so a kernel whose only driver is for
that, an older one for instance, still binds the device and drives what the two have in common.

\xpart{e} Cell 0 is 0, an SPI; cell 1 is 45, \code{0x2D}, the number among the SPIs; cell 2 is 1,
rising edge. The GIC's own numbering adds the SPI base of 32, so this is hardware interrupt
$45 + 32 = \mathbf{77}$, \code{0x4D}, edge-triggered on the rising edge.

\xpart{f} Yes: \code{acme,thing} is the node's second string, and one match is all that binding
needs. With only \code{acme,thing-v3} it does not bind, because that string is not in the node's
list, and \code{probe} is never called. A driver for a newer chip cannot claim an older one;
compatibility is declared by the device, not by the driver.

% ------------------------------------------------------------------------------------------
\solution[fillaccent3]{c10:ex:4}{What the helpers do}{Hand calculation}

\xpart{a} At the moment of the call it reads nothing from the tree. It takes the device's memory
resource 0, which the kernel built from \code{reg} entry 0, translated through every \code{ranges}
above it, when it created the device (\file{drivers/of/platform.c}:133); then it claims that region
and maps it. On success it returns the mapping, a \code{void __iomem *}. On failure it returns an
error packed into the pointer: \code{-EINVAL} if there is no such memory resource, \code{-EBUSY} if
the region is already claimed, \code{-ENOMEM} if the mapping fails
(\code{__devm_ioremap_resource}, \file{lib/devres.c}:135--165). It never returns \code{NULL}.

\xpart{b} Entry 0 of \code{interrupts}, with \code{interrupt-parent} to find the controller, or
\code{interrupts-extended}, which names both. It has that controller's IRQ domain translate the
specifier and create the mapping, and returns the Linux IRQ number, which is positive. On failure
it returns a negative error: \code{-EPROBE_DEFER} if the controller has not registered its domain
yet (\file{drivers/of/irq.c}:470), \code{-ENXIO} if there is no such interrupt, and \code{-EINVAL}
rather than 0 if the mapping produced 0 (\file{drivers/base/platform.c}:234--238). It never returns
0.

\xpart{c} The property \code{fifo-depth}, as one big-endian cell, which it stores in \code{v} in the
processor's own byte order. It returns 0 on success. On failure it returns \code{-EINVAL} if the
property does not exist, \code{-ENODATA} if it has no value, and \code{-EOVERFLOW} if it is too
short to hold a cell (\file{include/linux/of.h}:1350), and leaves \code{v} untouched, which is what
makes the default idiom of \secref{c10:app:b6} work.

\xpart{d} \code{platform_get_irq} returns a negative error code as an \code{int}, and a positive
number on success: test \code{irq < 0}, and return \code{irq} if it is.
\code{devm_platform_ioremap_resource} packs its error into the pointer: test \code{IS_ERR(regs)},
and return \code{PTR_ERR(regs)} if it is. \code{of_property_read_u32} returns zero for success and
delivers the value through its third argument: any non-zero return is the failure, which for an
optional property means ``use the default'' rather than a reason to give up.

\xpart{e} It works on every board where the call succeeds, because the number is positive and the
check passes, and that includes every board the driver was tested on. It silently does the wrong
thing on every board where the call fails, because \code{platform_get_irq} never returns 0 and the
check never fires. The negative error goes on to \code{devm_request_irq} as an \code{unsigned int},
a huge number with no descriptor behind it, and the request fails with \code{-EINVAL}
(\file{kernel/irq/manage.c}:2178--2180). The worst case is a board whose interrupt controller
probes after this driver: \code{platform_get_irq} returns \code{-EPROBE_DEFER}, which it does not
print, the driver turns it into \code{-EINVAL}, the core reports that the probe failed with error
-22, and the device is never tried again.

\xpart{f} From \secref{c6:app:b7}: \code{request_mem_region} and \code{ioremap}, the two
acquisitions, and \code{iounmap} and \code{release_mem_region}, the two releases. With \code{devm_}
the releases happen at unbind, so they disappear from the error path and from \code{remove} alike.

% ------------------------------------------------------------------------------------------
\solution[fillaccent4]{c10:ex:5}{What belongs in a device tree}{Design}

\xpart{a} The device tree, as \code{reg}. Where the block sits is a fact about the SoC and the board,
different on every design that uses the chip, so the driver must not contain it.

\xpart{b} The device tree, as \code{interrupts} and \code{interrupt-parent}. Which line of which
controller the device is wired to is again a fact about the board.

\xpart{c} A module parameter, or nothing at all. The depth of a software FIFO trades memory against
the risk of dropping data: a policy of this system's software, not a fact about the hardware, and
the same board may want different depths for different uses. A constant in the driver is right
until somebody needs to tune it, and a module parameter is right then. The \code{fifo-depth}
property that \secref{c10:app:b6} reads is legitimate only if it describes the device's own
hardware FIFO, as \code{QA_DEV_FIFO_DEPTH}, 16, does for \code{qa-dev}: that is hardware
description.

\xpart{d} Neither, normally. How many channels a chip has follows from which chip it is, and
\code{compatible} already says that; the driver looks the number up from the match, typically in
the \code{.data} of its \code{of_device_id} entry. A property is defensible only when the number
depends on the board rather than on the chip, for instance on which inputs are wired to anything,
and bindings usually describe that with a child node for each connected channel. That is what an
answer of ``the device tree'' has to argue.

\xpart{e} Not the device tree, which describes hardware. A module parameter at most, and on this
kernel not even that: \code{dev_dbg} with dynamic debug, which \file{kernel/qa.config} enables,
switches a driver's messages on and off at run time, call site by call site, with no parameter at
all (\secref{c4:app:a7}).

\xpart{f} The device tree: which crystal is fitted is a fact about the board. It is best described
as a clock rather than as a bare number: a \code{fixed-clock} node carrying \code{clock-frequency},
referenced from the device's \code{clocks} property, so that the driver asks the clock framework
for the rate (Chapter~\ref{c11:ch}).

\xpart{g} Neither. What a node in \file{/dev} is called is userspace's decision: the kernel's
subsystem gives it a default name, and udev or mdev rules rename it or add links to suit the
system. It is not a fact about the hardware, so it cannot be in the tree, and a parameter would move
a naming policy into the driver. The nearest thing a device tree legitimately carries is an alias
such as \code{serial0}, which numbers an instance and names no file.

\xpart{h} The device tree describes the hardware, what would be true of the board under any
operating system and any driver; everything that is a choice about how software uses that hardware
belongs somewhere else.

\xpart{i} Anything put in a device tree has to be supported for as long as a DTB containing it
exists. The DTB may come from firmware that is never updated, so every later kernel must accept it:
a property cannot be renamed, removed or given a new meaning, and a mistake is kept. The binding is
also shared with every other consumer of the tree, the bootloader and other operating systems,
which is why it is reviewed as hardware description. A module parameter is set on a command line or
in a configuration file that ships in the same system image as the module that reads it, and
changes with it: it is still an interface, visible under \file{/sys/module}, but the people who can
change it are the people who ship the kernel. The device tree often belongs to somebody else
entirely.

% ------------------------------------------------------------------------------------------
\solution[fillaccent4]{c10:ex:8}{Autoloading}{Design}

\xpart{a} The four steps:
\begin{enumerate}
  \item \code{MODULE_DEVICE_TABLE} has modpost turn the match table into alias strings in the
    \code{.ko}'s \code{.modinfo}.
  \item \code{depmod} collects the aliases of every installed module into \file{modules.alias}.
  \item The device's \code{uevent} carries a \code{MODALIAS=} string built from its node.
  \item udev receives the uevent and runs \code{modprobe} with that string, which finds the module
    through \file{modules.alias}.
\end{enumerate}
Step 4 is the one missing: there is no udev, the init script starts no \code{mdev}, and the kernel
is built without a uevent helper (\code{CONFIG_UEVENT_HELPER} is not set). For your own module step
2 is missing as well. It is loaded from \file{/lab} rather than installed under \file{/lib/modules},
so the \file{modules.alias} on the target, which \file{ci/rootfs.sh} copies across with the kernel's
own modules, has never heard of it.

\xpart{b} Step 1, and through it step 2. Without it the \code{.ko} advertises no alias,
\code{depmod} has nothing to collect, and \code{modprobe} given the device's \code{MODALIAS} finds no
module, so the driver is never loaded automatically. Loaded by hand it works perfectly, because
binding uses the \code{of_match_table} itself (\code{platform_match},
\file{drivers/base/platform.c}:1345) whether or not it was exported, which is why the omission goes
unnoticed on a target that loads everything by hand.

\xpart{c} \code{OF_COMPATIBLE_0=qacademy,qa-dev-1.0} is the node's \code{compatible} property,
copied into the uevent by \code{of_device_uevent} (\file{drivers/of/device.c}:238--239). QEMU wrote
that property into the tree it generated, from \code{QA_DEV_DT_COMPATIBLE} in \file{qa-dev.h}, in
\code{add_qa_dev_fdt_node} in \file{tools/qa-dev/integrate.py}. The \code{MODALIAS} line,
\code{of:Nqa-devT(null)Cqacademy,qa-dev-1.0}, is built from the same node by \code{of_modalias}
(\file{drivers/of/module.c}:29--50): \code{N} and the node's name, \code{T} and its
\code{device_type}, which this node does not have, and \code{C} before each \code{compatible}
string.

\xpart{d} Two aliases, which modpost generates from the one entry in the table
(\file{scripts/mod/file2alias.c}:366--381): \code{of:N*T*Cqacademy,qa-dev-1.0} and
\code{of:N*T*Cqacademy,qa-dev-1.0C*}, the second of which also matches a node that lists further
\code{compatible} strings after this one. \code{depmod} would collect them into
\file{/lib/modules/6.12.30/modules.alias}, with an indexed copy in \file{modules.alias.bin}; here it
has not, because the module is not installed there.

\xpart{e} The steps are the same; what differs is where the identity comes from. A USB device
describes itself: when it is plugged in, the USB core reads its descriptors, vendor, product and
class, and the uevent carries a \code{MODALIAS=usb:v...p...} built from them, which the aliases of
drivers declaring \code{MODULE_DEVICE_TABLE(usb, ...)} match. A device soldered to the board
describes nothing: a block of registers at an address has no descriptors to read and no bus that
could ask for them. Its identity has to come from outside, and the device tree is where it comes
from: the node's \code{compatible} is what the kernel creates the device from and builds its
\code{MODALIAS} out of. Without the node there is no device, no uevent and nothing to match.

% ------------------------------------------------------------------------------------------
\solution[fillaccent]{c10:ex:9}{An address derived three ways}{Cross-check}

\xpart{a} The bytes, which \secref{c10:app:b9} prints, grouped four to a cell, most significant
first, with the parent's \code{ranges} alongside:

\begin{console}
reg         00 00 00 00  00 00 10 00
            0x00000000   0x00001000
ranges      00 00 00 00  00 00 00 00  0c 00 00 00  02 00 00 00
            0x00000000   0x00000000   0x0c000000   0x02000000
interrupts  00 00 00 00  00 00 00 70  00 00 00 04
            0x00000000   0x00000070   0x00000004
\end{console}

\noindent The parent, \code{platform-bus@c000000}, declares one address cell and one size cell, so
\code{reg} is child address \code{0x0} and size \code{0x1000}. In \code{ranges} the child address is
one cell, \code{0x0}; the parent address is two, because the root declares
\code{\#address-cells = <2>}, so it is \code{0x00000000_0c000000}; and the size is one cell,
\code{0x02000000}. The device's window lies inside that range, so the physical address is
\code{0x0C000000 + (0x0 - 0x0) =} \textbf{\code{0x0C000000}}, the size \code{0x1000}, and the window
runs from \code{0x0C000000} to \code{0x0C000FFF}. The interrupt is type 0, an SPI; number
\code{0x70}, which is 112; flags 4, level high. The GIC hardware interrupt is
$112 + 32 = \mathbf{144}$, \code{0x90}.

\xpart{b} The name is \code{c000000.qa-dev}, and the number in it is the physical address,
\code{0xC000000}, not the \code{0x0} in \code{reg}. So the translation happened before any driver
was involved: when the kernel created the device at boot, which translated \code{reg} through
\code{ranges} and named the device from the result (\code{of_device_make_bus_id},
\file{drivers/of/device.c}:302).

\xpart{c} A correct driver prints \code{start} as \code{0xC000000} and \code{end} as
\code{0xC000FFF}, which is the last byte of the window rather than one past it
(\code{0xC000000 + 0x1000 - 1}), and the number \code{platform_get_irq} returned, which was 17 on the
course's target (\secref{c10:app:a6}).

\xpart{d} \textbf{They agree}, on \code{0x0C000000}. A hand figure of \code{0x0} took \code{reg} at
face value, or took the first parent cell, \code{0x00000000}, as the base. That cell is only the
high half of a two-cell address; the base is that cell and the next together,
\code{0x00000000_0c000000}.

\textbf{17 and 144} are both right because they are numbers in two different spaces
(\secref{c8:app:a1}). 144 is the GIC's hardware number, fixed by the wiring. 17 is the Linux IRQ
number, an index the kernel handed out when the GIC's IRQ domain first mapped hardware interrupt
144, and it depends on how many mappings were made before it. In \file{/proc/interrupts} the first
column, \code{17:}, is the Linux number; after the per-CPU counts come the controller,
\code{GIC-0}, its hardware number, \code{144}, and the trigger, \code{Level}.

\textbf{\code{-tx4}} reads each group of four bytes as one integer in the host's own byte order,
which on this target is little-endian, so it byte-swaps every cell:

\begin{console}
reg         00000000 00100000
ranges      00000000 00000000 0000000c 00000002
interrupts  00000000 70000000 04000000
\end{console}

\noindent \code{-tx1}, reassembled most significant byte first, is correct, because device tree
cells are big-endian on every machine (\secref{c10:app:b4}). The wrong answer is plausible because
four of the nine cells are zero and read the same either way, and the rest still look like numbers:
\code{0x0000000C} is as good an address as any to someone who has not seen the right one, and
\code{0x00100000} is a perfectly ordinary 1\,MiB window. Only \code{0x70000000} as an interrupt looks
absurd, and a reader who stops at the address never gets that far.

\xpart{e} \textbf{The address} is in the device tree QEMU generated when it started, the file
\code{dumpdtb} would have written (\secref{c10:app:b8}), which the kernel shows under
\file{/sys/firmware/devicetree/base}: node \code{/platform-bus@c000000/qa-dev@0}, property
\code{reg}, translated by the parent's \code{ranges}. The number itself is QEMU's choice: its
\code{virt} machine puts the platform bus at \code{0x0C000000} (\file{hw/arm/virt.c}:187), and the
first device placed on it goes at offset 0.

\textbf{The second device.} With a private structure allocated for each device, each \code{probe}
has its own registers, interrupt and counters, so the two \code{interrupts} attributes count two
different devices, as the two lines of \file{/proc/interrupts} in the exercise do, 426 and 420.
With a file-scope \code{static}, the second \code{probe} overwrites the first one's state: both
attributes read the same numbers, the handler for interrupt 17 reads and acknowledges the second
device's registers, so the first device's level-triggered line is never cleared, and unbinding
either device unmaps registers the other still uses. The Chapter~\ref{c6:ch} driver could not have
coped at all: it maps the one address its \code{base=} parameter names, a module can be loaded only
once, and so it drives exactly one device, at an address somebody had to type.

\textbf{\code{0xC001000}.} In Exercise~\ref{c6:ex:9} nothing was there: it is one page past the only
device, inside the platform bus's window but backed by nothing, and reading it took the machine
down. Now QEMU has placed a second \code{qa-dev} there, in the first free page-aligned slot on its
platform bus (\file{hw/core/platform-bus.c}), with the next free interrupt, SPI 113, which is GIC
145; and the same address is a working device. An address alone tells a driver nothing about what
is at it, or whether anything is; only the description can.

\textbf{A different SoC} needs nothing changed in the driver. That board's device tree supplies a
node with the same \code{compatible} and its own \code{reg}, \code{ranges} and \code{interrupts},
and the same source is rebuilt for that board's kernel. The driver changes only if the block itself
changes, and then by gaining an entry in its match table.

\xpart{f} The device tree is the board's description, written by whoever designed the board and
supplied with it; the driver is shared by every board that carries the chip. An address in the
driver is right on one board and wrong on all the others, and on those it claims and reads whatever
lives there, which Exercise~\ref{c6:ex:9} showed can take the machine down. It also hides the fault
it was meant to cover: a wrong tree should fail loudly, where its owner will see it and fix it, not
be papered over by a second copy of the same fact that will drift from the first. And the fallback
cannot do what it promises anyway. With no node there is no platform device and \code{probe} is
never called, so the only case in which the fallback could run is one in which the tree did supply
an address, and then the tree's address is the one to believe.
